\hypertarget{error-handling-robustness}{%
\chapter{ERROR HANDLING \& ROBUSTNESS}\label{error-handling-robustness}}

\hypertarget{error-handling-debugging-c98}{%
\chapter{ERROR HANDLING \& DEBUGGING
(C++98)}\label{error-handling-debugging-c98}}

\hypertarget{assertions}{%
\section{1. ASSERTIONS}\label{assertions}}

Use \passthrough{\lstinline!assert!} to catch programming logic errors
during development.

\begin{lstlisting}[language={C++}]
#include <iostream>
#include <cassert>

int divide(int a, int b) {
    assert(b != 0);  // Terminates program if b is 0
    return a / b;
}

int main() {
    std::cout << divide(10, 2) << "\n";
    return 0;
}
\end{lstlisting}

\begin{center}\rule{0.5\linewidth}{0.5pt}\end{center}

\hypertarget{exceptions}{%
\section{2. EXCEPTIONS}\label{exceptions}}

C++ uses exceptions to handle runtime errors gracefully.

\hypertarget{basic-try-catch}{%
\subsection{2.1 Basic Try-Catch}\label{basic-try-catch}}

\begin{lstlisting}[language={C++}]
#include <iostream>

int safe_divide(int a, int b) {
    if (b == 0) {
        throw "Division by zero condition!";
    }
    return a / b;
}

int main() {
    try {
        int result = safe_divide(10, 0);
        std::cout << result << "\n";
    } catch (const char* msg) {
        std::cerr << "Error: " << msg << "\n";
    }
    return 0;
}
\end{lstlisting}

\hypertarget{catching-multiple-types}{%
\subsection{2.2 Catching Multiple Types}\label{catching-multiple-types}}

\begin{lstlisting}[language={C++}]
try {
    // code...
} catch (int e) {
    std::cerr << "Integer exception: " << e << "\n";
} catch (const char* e) {
    std::cerr << "String exception: " << e << "\n";
} catch (...) {
    std::cerr << "Unknown exception caught!\n";
}
\end{lstlisting}

\hypertarget{standard-exceptions}{%
\subsection{2.3 Standard Exceptions}\label{standard-exceptions}}

Inherit from \passthrough{\lstinline!std::exception!} for custom
exceptions.

\begin{lstlisting}[language={C++}]
#include <iostream>
#include <exception>
#include <stdexcept> // runtime_error, logic_error

class MyError : public std::exception {
public:
    virtual const char* what() const throw() {
        return "My Custom Error";
    }
};

void test() {
    throw std::runtime_error("Standard Runtime Error");
}

int main() {
    try {
        test();
    } catch (const std::exception& e) {
        std::cerr << "Caught: " << e.what() << "\n";
    }
    return 0;
}
\end{lstlisting}

\hypertarget{exception-specifications-c98}{%
\subsection{2.4 Exception Specifications
(C++98)}\label{exception-specifications-c98}}

In C++98, you can specify what a function might throw. (Note: Deprecated
in C++11, but valid here).

\begin{lstlisting}[language={C++}]
// Can throw only int
void func() throw(int) {
    throw 42;
}

// Cannot throw anything
void no_throw() throw() {
    // logic...
}
\end{lstlisting}

\hypertarget{stack-unwinding-raii}{%
\subsection{2.5 Stack Unwinding \& RAII}\label{stack-unwinding-raii}}

When an exception is thrown, the stack is unwound, identifying
destructors for all local objects.

\begin{lstlisting}[language={C++}]
class Resource {
public:
    Resource() { std::cout << "Acquired\n"; }
    ~Resource() { std::cout << "Released\n"; }
};

void risky() {
    Resource r; // Destructor guaranteed to run
    throw 1;
}

int main() {
    try {
        risky();
    } catch (...) {
        std::cout << "Caught\n";
    }
    return 0;
}
// Output: Acquired -> Released -> Caught
\end{lstlisting}

\begin{center}\rule{0.5\linewidth}{0.5pt}\end{center}

\hypertarget{debugging-strategies}{%
\section{3. DEBUGGING STRATEGIES}\label{debugging-strategies}}

\hypertarget{debug-macros-c98-style}{%
\subsection{3.1 Debug Macros (C++98
Style)}\label{debug-macros-c98-style}}

\begin{lstlisting}[language={C++}]
#include <iostream>

#ifdef DEBUG
    #define LOG(x) std::cout << "DEBUG: " << x << "\n"
#else
    #define LOG(x) 
#endif

int main() {
    int x = 10;
    LOG("x is " << x);
    return 0;
}
\end{lstlisting}
